\documentclass[nomath]{harryopo-paper}

\title{图像识别研究报告}
\author{研究员}
\date{2026年08月08日}


\begin{document}

\maketitle

\section*{一、引言}
\addcontentsline{toc}{section}{一、引言}

本文研究\fzht{深度学习}在图像识别中的应用。

\subsection*{1.1 方法}
\addcontentsline{toc}{subsection}{1.1 方法}

实验配置如下表：

\begin{table}[htbp]
  \centering
  \small
  \begin{tabularx}{\textwidth}{>{\raggedright\arraybackslash}X >{\raggedright\arraybackslash}X >{\raggedright\arraybackslash}X}
    \toprule
    \fzht{模型} & \fzht{准确率} & \fzht{速度} \\
    \midrule
    ResNet-50 & 94.2\% & 1.2s \\
    EfficientNet & 95.1\% & 0.8s \\
    \bottomrule
  \end{tabularx}
  \caption{表1}
\end{table}

消融实验：

\begin{table}[htbp]
  \centering
  \small
  \begin{tabular}{p{0.317\textwidth} p{0.317\textwidth} p{0.317\textwidth}}
    \toprule
    \multicolumn{2}{c}{\fzht{配置}} & \fzht{结果} \\
    \midrule
    baseline & 90.5\% & 无增强 \\
    full & 96.8\% & 完整方案 \\
    \bottomrule
  \end{tabular}
  \caption{表2}
\end{table}

多数据集：

\begin{table}[htbp]
  \centering
  \small
  \begin{tabular}{p{0.475\textwidth} p{0.475\textwidth}}
    \toprule
    \fzht{模型} & \fzht{准确率} \\
    \midrule
    \multirow{2}{*}{Ours} & 98.5\% \\
     & 85.3\% \\
    \multirow{2}{*}{Base} & 95.2\% \\
     & 78.9\% \\
    \bottomrule
  \end{tabular}
  \caption{表3}
\end{table}

\section*{二、结论}
\addcontentsline{toc}{section}{二、结论}

实验证明所提方法有效。

\end{document}